\begin{quote}
\small
\justifying
It is widely believed that rich classical control, aggressive quantum optimization, and formal reasoning cannot coexist inside a single lightweight quantum programming system: one must either embed quantum operations inside a host language and sacrifice compiler-level semantic control, or restrict the language to static circuit fragments and sacrifice expressive hybrid programs. In this work we challenge this assumption and demonstrate that Eigen 2.3 "Helios", a runtime-first hybrid classical-quantum language implemented primarily in Python with optional Rust acceleration, occupies a practically useful middle point. Eigen exposes a full classical virtual machine with recursion, structured data, exceptions, arrays, maps, and parallel task constructs, while also lowering quantum regions to a dependency-aware DAG intermediate representation, applying local circuit rewrites, estimating resources, exporting to multiple backends, and invoking exact or ZX-based equivalence checking.

We present the first publication-style technical reconstruction of the system directly from the repository artifact, language specification, compiler/runtime implementation, benchmark suite, and observed command-line executions. The paper contributes four results. First, we formalize Eigen as a two-path architecture combining a stack bytecode VM and a topological EQIR runtime, and we derive invariants under which local quantum rewrites preserve semantics. Second, we characterize the simulator stack, including dense, sparse, matrix-product-state, GPU, and native Rust execution paths, and we show how the design cleanly separates hybrid expressivity from hardware-target portability. Third, using the shipped eight-program benchmark corpus, we report observed total end-to-end runtimes ranging from 0.833 ms to 573.935 ms on the evaluation host, with seven of eight quantum and hybrid programs completing in under 1.5 ms and the recursive `fibonacci.eig` workload surfacing the true cost of runtime recursion at 234 ms execution time rather than hiding it in a host interpreter. Fourth, we document both strengths and weaknesses: exact equivalence proofs for small circuits, reproducibility metadata emitted at commit `af2c9217`, but also a profiler crash on benchmarked programs, a recursion failure in resource estimation, and several specification-implementation mismatches.

Our central conclusion is that runtime-first design is not an anti-compiler position. Rather, in Eigen it functions as an organizing principle: dynamic classical semantics remain first-class, while static quantum subproblems are selectively extracted, optimized, verified, and exported.

\end{quote}
